\section{Every original unit forces a finite singularity in the fixed old conductor spaces}
The global existence result ANP2 has a direct consequence for the original
linear conductor itself, without selecting its moments independently.
Fix any original $0<\delta<1/2$, $\gamma>2$ and any nonzero original
unit $a=\alpha+i\chi$. Define the original entire function
\[
 G(z)=\mu_0(z)=E_A(0,z)=\prod_{j=1}^4F_j(z).
 \tag{GUC1}
\]
Every $F_j$ is nonconstant by ANP2, hence is not identically zero.
The identity theorem then implies that their product is not identically
zero. ANP2 supplies a nonzero finite zero $z_0$ of each chosen factor,
so in particular $G(z_0)=0$. Its multiplicity is a finite positive
integer $\nu$; write its complete local factorization as
\[
 G(z)=(z-z_0)^\nu h(z),\qquad
 h(z_0)=\frac{G^{(\nu)}(z_0)}{\nu!}\ne0.
 \tag{GUC2}
\]
This follows from the first nonzero Taylor coefficient of the nonzero
entire function. Its order equals the sum of the vanishing orders of
the four actual factors at $z_0$, with zero contribution from each
nonvanishing factor. No common factor or multiplicity is removed.

Retain every original Gamma measure, both centres and
$\rho_n=\sqrt{n!/(s/2)_n}$, with $s\in\{1,a_{\rm amp}\}$.
The fixed old $v=0$ conductor in degree $D$ is precisely FPZ15,
with $g(t,z)=e^{-4\omega(t)}E_A(\omega(t),z)$ and
$\omega(t)=-i\arctan t$. Its diagonal is $g(0,z)=G(z)$.
All entries are entire in $z$, since each required coefficient
is a finite derivative in $\xi$ of the original factors. Therefore
\[
 \begin{split}
 B_{D,s}(z)&=\operatorname{diag}(\rho_0,\ldots,\rho_D)^{-1}
              T_D(g(z))\operatorname{diag}(\rho_0,\ldots,\rho_D),\\
 \det B_{D,s}(z)&=G(z)^{D+1},\qquad
 \operatorname{Spec}B_{D,s}(z)=\{G(z)\text{ repeated }D+1\text{ times}\}.
 \end{split}\tag{GUC3}
\]
Triangularity proves both identities. In particular the original
constant polynomial is killed at $z_0$ and the full matrix is singular.
For $z\ne z_0$ sufficiently close to $z_0$, GUC2 makes the matrix
invertible. Its inverse remains triangular with diagonal $G(z)^{-1}$.
In the ordered exterior basis every diagonal of its $r$th exterior
power is therefore $G(z)^{-r}$. The operator norm is at least the
modulus of any matrix entry in the two orthonormal exterior frames:
apply it to the relevant source basis vector and project onto the
relevant target basis vector. No identification of physical centres
is required for this coordinate argument.
This proves in the exact original orthonormal Gamma coordinates
\[
 \begin{gathered}
 \|\bigwedge^r B_{D,s}(z)^{-1}\|\ge |G(z)|^{-r},
               \qquad 1\le r\le D+1,\\
 \liminf_{z\to z_0}|z-z_0|^{\nu r}
       \|\bigwedge^r B_{D,s}(z)^{-1}\|
            \ge |h(z_0)|^{-r}>0,\\
 \lim_{z\to z_0}|z-z_0|^{\nu(D+1)}
       \|\bigwedge^{D+1}B_{D,s}(z)^{-1}\|
            =|h(z_0)|^{-(D+1)}>0.
 \end{gathered}\tag{GUC4}
\]
The last equality uses the one-dimensional top exterior space and
the exact determinant GUC3. Thus every inverse exterior rank diverges
at a finite nonzero original period for every original quartet and unit.
The lower bound does not assert that the smaller exterior ranks have
exactly the lower exponent $\nu r$; UPP18 calculates their stronger
exact exponent at the certified order-one points.

There is no hidden assumption about symbol order at the guaranteed
zero. If the entire function $E_A(\xi,z_0)$ is nonzero and has order
$v\ge1$ at $\xi=0$, then $g(t,z_0)$ has the same order, since
$\omega'(0)=-i\ne0$ and $e^{-4\omega(0)}=1$.
Writing $g=t^v\widetilde g$ with $\widetilde g(0)\ne0$ shows that
$T_D(g)$ is a $v$-step shift times an invertible triangular matrix.
Consequently its rank is $\max(D+1-v,0)$ and its kernel is the
original $\mathcal P_{\min(v-1,D)}$.
If instead $E_A(\xi,z_0)$ is identically zero, every coefficient
of $g$ vanishes and $B_{D,s}(z_0)=0$. GUC3--4 hold in that case
as well, because they depend on the nonzero entire function $G(z)$
of the period parameter, not on a finite symbol order at $z_0$.
This exhausts both cases. Neither orbit admissibility nor a finite
symbol-order hypothesis was used to infer the fixed-space singularity.

Finally put $u_0=1/z_0$, retaining the original finite nonzero period.
The exact identity
$z-z_0=-(u-u_0)/(uu_0)$ gives
\[
 \begin{split}
 \liminf_{u\to u_0}|u-u_0|^{\nu r}
       \|\bigwedge^r B_{D,s}(u^{-1})^{-1}\|
       &\ge |u_0|^{2\nu r}|h(z_0)|^{-r},\\
 \lim_{u\to u_0}|u-u_0|^{\nu(D+1)}
       \|\bigwedge^{D+1}B_{D,s}(u^{-1})^{-1}\|
       &=|u_0|^{2\nu(D+1)}|h(z_0)|^{-(D+1)}.
 \end{split}\tag{GUC5}
\]
These statements locate a necessary obstruction to extending an
inverse-exterior estimate across the whole complex period family.
They do not place the guaranteed period inside any separately
restricted arithmetic period domain. UPP3--20 provides explicit
locations, actual symbol order, five-orbit admissibility and the
corrected ES return on its certified geometric region; the global
conclusion here is the proved singularity and divergence GUC1--5.
The complete original coefficient and norm sources remain
\cite{PCLOriginal,WCF}; ANP1--32 supplies the new existence input.
